% \documentclass{article}

% % Language setting
% % Replace `english' with e.g. `spanish' to change the document language
% \usepackage[english, russian]{babel}

% % Set page size and margins
% % Replace `letterpaper' with `a4paper' for UK/EU standard size
% \usepackage[letterpaper,top=2cm,bottom=2cm,left=3cm,right=3cm,marginparwidth=1.75cm]{geometry}

% % Useful packages
% \usepackage{amsmath}
% \usepackage{graphicx}
% \usepackage[colorlinks=true, allcolors=blue]{hyperref}

\documentclass[12pt,a4paper]{article}
\usepackage[T2A]{fontenc}
\usepackage[utf8]{inputenc}
\usepackage[russian]{babel}
\usepackage{geometry}
\geometry{top=2.5cm, bottom=2.5cm, left=2cm, right=2cm}
\usepackage{amsmath, amssymb, amsfonts}
\usepackage{tcolorbox}
\tcbuselibrary{breakable}
\usepackage{xcolor}
\usepackage{hyperref}

\definecolor{examplebg}{HTML}{F5F5F5}
\definecolor{exampleframe}{HTML}{CCCCCC}

\newtcolorbox{examplebox}[1][]{%
  breakable, 
  colback=examplebg, 
  colframe=exampleframe, 
  title=Пример, 
  fonttitle=\bfseries, 
  #1
}

\begin{document}

\large

\section{Вариант 2}

\subsection{Задача 1}

\textbf{Описание:}

Вы играете на выживании и провалились на дно глубокого ущелья. Расстояние от дна до края поверхности — ровно 30 блоков. 

Вы начинаете строить столб из булыжника прямо под собой (делать «пиллар»). Днём вы находитесь в безопасности и успеваете построить 6 блоков вверх. 
Но с наступлением ночи в небе появляются опасные Фантомы. Чтобы укрыться от них до рассвета, вам приходится сломать 2 верхних блока своего же столба и спрятаться внутри образовавшегося углубления (из-за чего вы опускаетесь на 2 блока вниз). 

Утром Фантомы сгорают на солнце, и вы продолжаете строить столб с той отметки, где переждали ночь. Этот цикл повторяется каждый день.

Через сколько полных дней (считая от начала первого дня) ваш столб достигнет края ущелья, и вы сможете благополучно выбраться на поверхность? 

\textit{В бланк ответов впишите только одно число. Никаких пояснений к этой задаче писать не нужно.}

\newpage

\subsection{Задача 2}

\textbf{Описание:}

Вы пытаетесь открыть секретную дверь в сокровищницу древнего Храма в джунглях. Замок представляет собой магический круг из 16 сундуков (пронумерованных по кругу от 1 до 16). Изначально все шестнадцать сундуков абсолютно пусты. 

У вас есть специальный механизм из раздатчиков и воронок: при нажатии на кнопку он кладет ровно по 1 алмазу в любые два соседних сундука в круге одновременно. Нажимать на кнопку можно сколько угодно раз, подключая механизм к любым парам соседних сундуков.

Чтобы массивная железная дверь открылась, сундуки должны образовать «пароль»: в первом сундуке должен лежать 1 алмаз, во втором — 2 алмаза, в третьем — 3, и так далее вплоть до 16 алмазов в шестнадцатом сундуке. 
Сможете ли вы открыть дверь с помощью этого механизма? Обоснуйте свой ответ.

\textit{В бланк ответов запишите "Да" или "Нет", а затем подробно распишите ваше доказательство.}


\newpage

\subsection{Задача 3}

\textbf{Описание:}

На сервере проходит турнир за звание лучшего игрока. В финале на PvP-арене встречаются два игрока: Стив и Алекс. В центре арены стоит сундук, в котором лежит ровно 100 алмазов. 

Правила дуэли:

1. Стив и Алекс делают ходы по очереди. Стив ходит первым.

2. В свой ход игрок должен забрать из сундука 1, 2, 3 или 4 алмаза. Взять 0 алмазов (пропустить ход) нельзя.

3. Победителем турнира признается тот, кто заберет самый последний алмаз из сундука. Он же забирает себе весь выигрыш.

Оба дуэлянта супер-умные, отлично знают математику и не совершают ошибок. Кто победит в этой дуэли: Стив или Алекс? И самое главное — как именно он должен действовать, чтобы гарантированно победить независимо от ходов соперника?

\textit{В бланк ответов запишите, кто победит (Стив или Алекс), а затем опишите его выигрышную стратегию.}


\newpage

\subsection*{Задача 4}

\textbf{Описание:}

Вы пытаетесь пройти древнее испытание в мире Minecraft, чтобы добраться до сокровищницы. Перед вами коридор из $13$ последовательных редстоун-ворот.

У каждых ворот есть редстоун-лампа, которая может находиться в одном из двух состояний: «ВЫКЛ» или «ВКЛ». Изначально на всех 13 воротах лампы находятся в состоянии «ВЫКЛ».

Когда игрок подходит к любым воротам, происходит следующее:
\begin{enumerate}
    \item Игрок смотрит на текущее состояние лампы.
    \item Лампа мгновенно переключается в противоположное состояние (если была «ВЫКЛ», становится «ВКЛ», и наоборот).
    \item После этого механизм решает судьбу игрока в зависимости от того, в каком состоянии лампа была изначально:
    \begin{itemize}
        \item Если лампа БЫЛА «ВЫКЛ», срабатывает ловушка, и игрок немедленно телепортируется обратно к началу, к 1-м воротам. Приходится начинать попытку заново.
        \item Если лампа БЫЛА «ВКЛ», ворота открываются, и игрок сразу проходит к следующим воротам (или попадает в сокровищницу, если это были последние, 13-е ворота).
    \end{itemize}
\end{enumerate}

Каждый проход с 1-х ворот (включая самый первый) считается одной попыткой.

Какое минимальное количество попыток потребуется игроку, чтобы успешно пройти все 13 редстоун-ворот и попасть в сокровищницу?

\textit{В бланк ответов запишите одно число — суммарное количество попыток, а затем приведите обоснование минимальности.}

\newpage

\subsection*{Задача 5}

\textbf{Описание:}

В мире Minecraft есть прямой туннель длиной ровно $100$ блоков.
В этот туннель одновременно запустили $8$ вагонеток. Туннель настолько узкий, что вагонетки не могут разъехаться.

Каждая вагонетка движется со скоростью $1$ блок в тик. Изначально некоторые из них едут вправо (к отметке 100), а другие — влево (к отметке 0).
Если две вагонетки сталкиваются лоб в лоб, они мгновенно отскакивают друг от друга и меняют направление на противоположное, не теряя скорости.
Как только вагонетка достигает отметки $0$ или $100$, она навсегда выезжает из туннеля.

Игрок зафиксировал стартовые координаты всех 8 вагонеток (слева направо) и их направления:
\begin{itemize}
    \item Вагонетка 1: координата $10$, едет вправо (R)
    \item Вагонетка 2: координата $25$, едет вправо (R)
    \item Вагонетка 3: координата $30$, едет влево (L)
    \item Вагонетка 4: координата $45$, едет вправо (R)
    \item Вагонетка 5: координата $50$, едет влево (L)
    \item Вагонетка 6: координата $70$, едет влево (L)
    \item Вагонетка 7: координата $85$, едет вправо (R)
    \item Вагонетка 8: координата $95$, едет влево (L)
\end{itemize}

\textbf{Требуется вычислить:}
\begin{enumerate}
    \item Какое суммарное количество столкновений произойдет между всеми вагонетками, пока туннель полностью не опустеет?
    \item Через сколько тиков после старта туннель покинет самая последняя вагонетка?
\end{enumerate}

\textit{В бланк ответов запишите два числа: количество столкновений и время (в тиках). Запишите ваше решение.}

\end{document}
